\section{Radon--Nikodym 导数把``相对于谁''说清楚}
\label{sec:12-Real-Analysis-Support/05-Measure-Integration-and-Conditional-Expectation/04}

把接口延迟的单位从秒改成毫秒，重新拟合，密度函数的数值全部变成原来的千分之一。数据没换，分布没换，落在任何一段区间里的概率一个都没变，可 \(p(x)\) 这条曲线整体塌了三个数量级。

同一件事还有另一副面孔：方差很小的正态分布，密度在中心处可以远大于 1。概率不许超过 1，密度却随便超。

两处不适都指向同一件事：\(p(x)\) 不是分布自己携带的数值。它是一个比值，分子是分布分配的质量，分母是某把尺子分配的长度。换尺子——换单位、换坐标、换参考测度——比值就跟着变。把分母藏起来不写，密度看着像绝对量，各种困惑就是这么来的。

\subsection{绝对连续是尺子能看见的前提}

设 \(\nu\) 与 \(\mu\) 定义在同一个可测空间上。若
\[
\mu(A)=0\ \Longrightarrow\ \nu(A)=0,
\]
就说 \(\nu\) 对 \(\mu\) 绝对连续，记作 \(\nu\ll\mu\)。这句话的意思是 \(\nu\) 不往 \(\mu\) 看不见的地方堆质量。

为什么它是写出密度的前提，一眼就能看出来。若能写成 \(\nu(A)=\int_A f\,d\mu\)，那么 \(\mu(A)=0\) 时右边必为零，于是 \(\nu(A)\) 也必须为零。反过来，这条必要条件也够用。

\begin{theorem}[Radon--Nikodym]
设 \(\mu,\nu\) 是同一可测空间上的 \(\sigma\)-有限测度且 \(\nu\ll\mu\)，则存在非负可测函数 \(f\)，使得对一切可测集 \(A\) 有 \(\nu(A)=\int_A f\,d\mu\)；\(f\) 在 \(\mu\)-几乎处处意义下唯一。
\end{theorem}

这个 \(f\) 记作 \(\dfrac{d\nu}{d\mu}\)，称为 \(\nu\) 相对于 \(\mu\) 的 Radon--Nikodym 导数。唯一性只到``几乎处处''为止：在一个 \(\mu\)-零测集上随便改动 \(f\) 的值，所有积分都不受影响，所以密度从来不是一条被唯一钉死的曲线，而是一个等价类。教科书上正态密度写成那个漂亮公式，只是等价类里最光滑的那个代表。

\(\sigma\)-有限指空间能写成可数个有限测度集合的并。Lebesgue 测度、计数测度、任何概率测度都满足，实际工作里这条几乎不构成限制，但在无穷维空间上它会失效——那里通常没有``Lebesgue 测度''可用，也就没有默认的尺子。

\subsection{质量函数与密度是同一个对象换了尺子}

上一节说到，同一个换元公式因参考测度不同而显出不同长相。现在可以把那句话讲完。

设 \(P\) 是可数集上的概率分布，取计数测度 \(\#\) 作尺子，则
\[
\frac{dP}{d\#}(x)=P(\set{x}),
\]
这就是概率质量函数。设 \(P\) 是 \(\R^{d}\) 上对 Lebesgue 测度 \(\lambda\) 绝对连续的分布，则
\[
p(x)=\frac{dP}{d\lambda}(x),\qquad P(A)=\int_A p(x)\,dx,
\]
这就是通常的概率密度。两者是同一个构造，差别只在分母。``离散用求和、连续用积分''不是两套理论，是同一套理论的两次选尺。

选尺这个动作也解释了开头那两处不适。换单位相当于把 \(\lambda\) 整体缩放，比值自然反向缩放；多元变量换坐标时冒出来的 Jacobian，同样是在折算参考体积元，不是在改概率。密度大于 1 也就没什么可奇怪的了：比值可以任意大，只要积起来是 1。

\subsection{有没有密度，取决于跟谁比}

\begin{center}
\begin{tikzpicture}[>=stealth,scale=1.2]
\begin{scope}
  \draw[->] (-0.4,0) -- (3.4,0);
  \draw[thick] plot[domain=-0.3:3.2,samples=60] (\x,{0.85*exp(-1.6*(\x-1.9)*(\x-1.9))});
  \draw[->,very thick] (0.35,0) -- (0.35,1.75);
  \node at (0.35,2.05) {\footnotesize \(?\)};
  \node at (1.5,-0.5) {\footnotesize 尺子取 \(\lambda\)};
  \node at (1.5,-0.95) {\footnotesize 点质量处比值无处安放};
\end{scope}
\begin{scope}[shift={(4.6,0)}]
  \draw[->] (-0.4,0) -- (3.4,0);
  \draw[thick] plot[domain=-0.3:3.2,samples=60] (\x,{0.85*exp(-1.6*(\x-1.9)*(\x-1.9))});
  \draw[very thick] (0.35,0) -- (0.35,0.9);
  \fill (0.35,0.9) circle (0.06);
  \node at (0.35,1.2) {\footnotesize \(\alpha\)};
  \node at (1.5,-0.5) {\footnotesize 尺子取 \(\delta_0+\lambda\)};
  \node at (1.5,-0.95) {\footnotesize 同一个分布处处有密度};
\end{scope}
\end{tikzpicture}
\end{center}

\emph{分布没有变，变的是分母；密度存在与否是分布与尺子之间的关系，不是分布单方面的属性。}

点质量 \(\delta_0\) 把全部概率压在一点上。因为 \(\lambda(\set{0})=0\) 而 \(\delta_0(\set{0})=1\)，绝对连续当场失败，不存在任何普通函数 \(p\) 使 \(\delta_0(A)=\int_A p\,d\lambda\)。物理文献里写的 Dirac ``delta 函数''是别的东西，按普通密度去积分会出错。

混合分布同样如此。取
\[
P=\alpha\,\delta_0+(1-\alpha)\,\mathcal N(0,1),\qquad 0<\alpha<1,
\]
它有一块质量钉在 0 上，其余摊开，相对 \(\lambda\) 没有密度。这类分布在实践里很常见：一个计费金额，大部分订单是 0 元，其余连续分布；一个等待时长，有一批请求命中缓存直接返回 0。把它硬套连续密度，或者硬当离散处理，都会丢掉一半。

换一把尺子，问题就消失了。令 \(\mu=\delta_0+\lambda\)，则 \(P\ll\mu\)，且
\[
\frac{dP}{d\mu}(x)=
\begin{cases}
\alpha, & x=0,\\[2pt]
(1-\alpha)\,\dfrac{1}{\sqrt{2\pi}}e^{-x^{2}/2}, & x\ne 0.
\end{cases}
\]
在 \(x=0\) 处这个值是相对于 \(\delta_0\) 的比值，在别处是相对于 \(\lambda\) 的比值，两段拼成一个函数，最大似然、Bayes 更新、KL 计算都可以照常做。所以``这个分布有没有密度''是个不完整的问句，补全了才有答案：相对于哪把尺子。

\subsection{换测度就是给样本重新配权}

线上样本按策略 \(Q\) 采集，你要估计策略 \(P\) 下的平均收益。若 \(P\ll Q\)，令 \(w=\dfrac{dP}{dQ}\)，则对可积的 \(g\)，
\[
\E_P[g]=\int g\,dP=\int g\,w\,dQ=\E_Q[g\,w].
\]
手上只有 \(Q\) 的样本，照样能估 \(P\) 下的期望，代价是每条样本乘一个权重。重要性采样、离策略评估、covariate shift 校正，都是这一行式子。

这行式子也把两个条件摆到了明面上。第一个是支撑覆盖：\(P\ll Q\) 不成立时，\(P\) 在 \(Q\) 从不访问的区域放了质量，那部分收益没有任何样本携带信息，再多同源数据也补不回来——这是可识别性的问题，不是样本量的问题。第二个是权重的可控性：即便绝对连续成立，\(w\) 也可能在一小片区域上大得离谱，少数几条样本主导整个估计，方差爆掉；极端情形下 \(\E_Q[w^{2}]=\infty\)，估计量连中心极限定理都用不上。第六节讲的一致可积性，量的正是这类失控。

回到开头的混合类型数据。样本 \(x\) 里既有连续的延迟又有离散的设备类型，权重 \(p(x)/q(x)\) 之所以仍然写得出来，是因为两个分布可以共用一把乘积尺子——连续分量用 Lebesgue、离散分量用计数——比值里的分母约掉了，剩下的东西与选哪把尺子无关。这个约分是重要性权重能在混合数据上照常工作的原因。

\subsection{似然比与 KL 散度都是这个导数的读法}

同一个参数族里的两个模型 \(P_{\theta_0},P_{\theta_1}\)，若被共同的参考测度 \(\mu\) 支配，密度写作 \(p_\theta=dP_\theta/d\mu\)，则似然比
\[
\frac{p_{\theta_1}(x)}{p_{\theta_0}(x)}=\frac{dP_{\theta_1}}{dP_{\theta_0}}(x)
\]
与 \(\mu\) 无关。这一点值得记住：单个密度依赖尺子，两个密度的比不依赖。假设检验、Bayes 因子、判别器输出，比较的都是这个与尺子无关的量，所以它们不会因为你换了单位而改变结论。

KL 散度的一般定义正是对数 Radon--Nikodym 导数的期望：
\[
\KL(P\Vert Q)=\int\log\frac{dP}{dQ}\,dP=\E_P\Bigl[\log\frac{dP}{dQ}\Bigr].
\]
写成这个样子，离散情形的求和与连续情形的积分是同一条式子，见 \hyperref[sec:07-Information-Theory/03-Divergences/01]{``错误模型的代价：从交叉熵到 KL 散度''}。定义里内建了 \(P\ll Q\) 的要求，一旦有一块 \(Q\) 判为不可能而 \(P\) 仍会产出的区域，散度就是 \(\infty\)。这不是定义的粗糙，是 \(\KL\) 方向性的来源：拿一个把某些结果彻底排除的模型去解释真会发生的数据，代价无上限。

Bayes 更新也是同一件事的一次归一化。先验与似然若相对同一把尺子有密度，后验密度就是
\[
p(\theta\given x)=\frac{p(x\given\theta)\,p(\theta)}{\int p(x\given u)\,p(u)\,du}.
\]
这条公式熟悉到容易忘记它有前提。遇到离散连续混合的参数、函数空间上的先验、或者以零概率事件为条件，第一步不是套公式，是先确认支配测度存在——存在了公式照写，不存在就得换别的构造。

下一节要处理的正是``不存在也得办''的那一类。以一个零概率事件为条件求期望，用初等公式除下去会得到 \(0/0\)；能救场的恰是本节这个定理——条件期望的存在性，就是 Radon--Nikodym 定理的一次应用。
